\section{The Intent Elaboration Layer}
\label{sec:intent_layer}

The layered model developed earlier in this article places a machine-checkable formal contract at the pivot between the ambiguous human world and the precise machine world. Everything below that pivot---synthesis of code from the contract and verification of code against it---can in principle be made sound, because both the agent and the verifier operate over the same formal object. Everything above the pivot is harder, because it must carry intent across the boundary between a human who reasons informally and a machine that reasons formally, and it must do so without quietly substituting the machine's reading of intent for the human's. This is the work of the \emph{intent elaboration layer}: the process by which natural-language (NL) intent becomes a formal contract that the human accepts as a faithful statement of what they meant. It is the hardest layer in the pipeline precisely because its correctness cannot be delegated to a solver. Z3 can prove that code satisfies a contract, but no solver can certify that the contract says what the human wanted; that judgment is irreducibly human, and the entire assurance argument of the pipeline rests on it. This section defines the central tension of the layer, develops the mechanisms that resolve it, gives the in-layer elaboration loop as an algorithm, and states the layer's named research contribution.

\subsection{The Ratification Gap}
\label{subsec:ratification_gap}

The dual role of the specification, established earlier, is that the contract is simultaneously the agent's \emph{input}---the precise intent it builds against---and its \emph{oracle}---the property its output is verified against. As an oracle the contract is load-bearing in a way that a test suite never is: verification covers all inputs, not a sample, so a contract that misstates intent will be discharged confidently against code that does the wrong thing. A wrong specification verified against is therefore worse than no specification at all, because it converts the absence of a counterexample into false assurance. It follows that the contract must be correct with respect to intent and must be owned by the human, since only the human is the authority on what was meant.

Here the layer encounters its defining obstacle, which we term the \emph{ratification gap}. The human supplied intent in natural language for a reason: they cannot, or will not, write formal contracts. Asking that same human to confirm formal text---to read a quantified JML postcondition and attest that it captures their intent---makes the human's authority a fiction. They will sign whatever they are shown, because they cannot evaluate it, and the signature certifies nothing. The naive elaboration loop, in which an LLM drafts a contract and the human is asked ``is this correct?'', thus produces the appearance of human ownership while transferring real authority to the model that drafted the text. Because that model is the same kind of system whose output the contract is meant to police, the loop is circular in the most damaging way: the artefact that is supposed to ground the pipeline's assurance is itself ungrounded. Resolving the ratification gap without abandoning either human ownership or formal precision is the problem the remaining mechanisms of this layer are designed to solve.

\subsection{Ratify Behaviour, Not Logic}
\label{subsec:ratify_behavior}

The first and most important mechanism is a discipline: \emph{ratify behaviour, not logic}. The human is never asked to confirm formal text. Instead the candidate contract is rendered into concrete, checkable behavioural claims that any competent reader of the problem domain can evaluate without reading a line of JML. A contract clause is exercised at representative and boundary inputs, and the resulting input--output pairs are presented as questions in the human's own terms: ``given an input of $-1$, this contract says the result is $0$; is that correct?''; ``given an empty list, this contract says an exception is thrown rather than zero returned; is that what you want?''. The human reasons about behaviour, which is what they understand, rather than about the predicate-logic encoding of behaviour, which they do not.

Each behavioural claim the human ratifies becomes a \emph{pinned oracle}: a concrete input--output (or input--exception) pair that the contract is henceforth required to satisfy, recorded independently of the formal text. Pinned oracles serve three functions. They are the operational definition of human ratification, replacing an unverifiable signature on formal text with a verifiable set of agreed cases. They constrain the formal contract directly, because any subsequent reformulation---whether by the elaborator or by a later compositional pass---must continue to satisfy every pinned oracle, and a reformulation that does not is rejected mechanically rather than silently accepted. And they are themselves machine-checkable: the pinned oracle ``$f(-1)=0$'' can be discharged against the formal contract without human involvement, turning ratification into a persistent, re-checkable artefact rather than a one-time conversation. Example-grounded ratification is, in effect, the elaboration-layer analogue of programming by example~\cite{gulwani2011flashfill}: concrete instances anchor an otherwise underdetermined formal object. The difference is that here the examples do not synthesize the contract---inference and the LLM draft do that---but \emph{ratify} it, supplying the human authority that synthesis alone cannot. Pinned oracles are, however, necessarily a finite sample, so they do not by themselves establish human authority over the contract's behaviour on \emph{every} input; what they establish is authority over a strategically chosen, catalogue- and adversarially-guided set of boundary cases, while the contract's universal force on unpinned inputs is borne by the formal generalisation and the adversarial search. This is a different posture from the sampled-oracle weakness criticised earlier: there the unsampled region is simply unverified, whereas here a finite sample \emph{anchors} a universal contract that is then verified over all inputs, so the residual risk is contract-misstatement on unpinned inputs---a narrower hazard than test-suite coverage gaps, though not an eliminated one.

\subsection{Socratic, Decision-Surfacing Elaboration}
\label{subsec:socratic}

Behavioural ratification answers how to confirm a contract; it does not by itself decide which questions to ask. Natural-language intent is radically underdetermined at exactly the points where defects concentrate, and a faithful elaborator must resolve those points to produce any contract at all. Rather than resolve them silently and present a finished contract, the elaborator practices \emph{Socratic}, decision-surfacing elaboration: it makes explicit the ambiguities it was forced to resolve and asks the human only about those. The premise is that the human's scarce attention should be spent on genuine decisions, not on confirming the obvious.

The decisions worth surfacing are not arbitrary; they fall into a small, recurring \emph{boundary-decision catalogue} that the elaborator walks for each method. \Cref{tab:boundary_catalogue} lists the categories. For each one the elaborator determines whether the NL intent settles the decision; where it does not, the unresolved decision is surfaced as a behavioural question in the sense of \cref{subsec:ratify_behavior}. This converts an open-ended ``is this contract right?'' interrogation, which the human cannot answer, into a bounded sequence of concrete adjudications, each of which the human can. It also makes the elaboration auditable: the catalogue is a checklist, and a contract can be reported as having every catalogue entry either resolved by stated intent, resolved by an inherited default (\cref{subsec:inherited_defaults}), or ratified by the human.

\begin{table}[t]
\centering
\caption{The boundary-decision catalogue. For each method, the elaborator determines whether NL intent settles each category; unresolved categories are surfaced to the human as behavioural questions.}
\label{tab:boundary_catalogue}
\begin{tabular}{@{}ll@{}}
\toprule
Category & Representative question \\
\midrule
Null inputs        & Reject \texttt{null} arguments, or treat as a valid value? \\
Empty inputs       & Empty collection/string: return identity, or error? \\
Overflow           & Saturate, wrap, or throw on arithmetic overflow? \\
Ordering           & Is output order specified, or unconstrained? \\
Duplicates         & Are duplicate elements removed, preserved, or rejected? \\
Concurrency        & Is the method required to be thread-safe? \\
Error vs.\ exception & Signal failure by return value (sentinel) or by exception? \\
\bottomrule
\end{tabular}
\end{table}

\subsection{Adversarial Elaboration and Vacuity Checks}
\label{subsec:adversarial}

Surfacing decisions and ratifying behaviour guard against a contract that omits intent, but they do not guard against a contract that is technically consistent with the stated intent yet too weak to be useful. An LLM asked to formalise ``sort the list'' may produce a postcondition asserting only that the output is a permutation of the input---true of the sorting it was asked for, but also satisfied by the identity function. The drafting model has, perhaps inadvertently, weakened the specification to one its own implementation can trivially meet. Because the same family of models drafts the contract and synthesizes the code, this failure mode is systematic rather than accidental, and behavioural ratification will not always catch it: a human who agrees that ``the output is a permutation of the input'' is a true statement may not notice that it is an \emph{incomplete} one.

The mechanism that addresses this is \emph{adversarial elaboration}, the elaboration-layer dual of counterexample-guided inductive synthesis (CEGIS)~\cite{solarlezama2008sketching}. Where CEGIS pits a synthesizer against a verifier that searches for inputs on which the candidate program fails, adversarial elaboration pits the contract drafter against a second, independent agent that searches for \emph{scenarios consistent with the NL intent but violating the proposed contract's evident purpose}---concretely, behaviours a reasonable reading of the intent would forbid that the drafted contract nonetheless permits. For the sorting example, the adversary exhibits the identity function as a witness that satisfies the permutation-only contract while plainly contradicting ``sort'', signaling that the postcondition must be strengthened with an orderedness clause. The adversary's witnesses are themselves behavioural and are fed back into the ratification dialogue, so spec weakening is surfaced as a concrete ``the contract currently permits this; did you intend to permit it?'' question.

Adversarial elaboration is complemented by automatic \emph{vacuity and non-triviality checks}, which detect the degenerate cases of weakening without an adversary in the loop. Two patterns are flagged: a postcondition that is logically equivalent to \lstinline|true|, which constrains nothing, and a precondition that is unsatisfiable, which makes the contract vacuously dischargeable because no input ever reaches it. This repurposes a signal already present in our prior inference work, where the inferrer emits an \lstinline|ensures true| clause as an honest \emph{punt}---an explicit admission that it could not infer a meaningful postcondition for a method~\cite{edmonds_article1}. In the elaboration layer the same syntactic signal is reinterpreted as a weakening alarm: an \lstinline|ensures true| obligation, or any clause provably equivalent to it, marks a method whose contract carries no behavioural content and must either be strengthened or be explicitly accepted as unconstrained. \Cref{lst:vacuity} illustrates the two patterns.

\begin{lstlisting}[style=code,caption={Vacuity patterns flagged by non-triviality checks: a postcondition equivalent to \texttt{true} (no constraint) and an unsatisfiable precondition (vacuously dischargeable).},label={lst:vacuity}]
//@ ensures true;                  // flagged: no behavioural content
public int normalise(int x) { ... }

//@ requires n > 0 && n < 0;       // flagged: unsatisfiable -> vacuous
//@ ensures \result == n;
public int echo(int n) { ... }
\end{lstlisting}

\subsection{The Independent Authority Constraint}
\label{subsec:independent_authority}

Adversarial elaboration already requires a second agent, and this reflects a constraint that governs the whole layer: the \emph{independent authority constraint}. The contract that code is verified against must not be authored by the same agent---the same model instance, and more importantly the same role---that writes the code. If a single agent drafts both the oracle and the implementation, then a misreading of intent corrupts both identically, the implementation satisfies the corrupted oracle, and verification certifies the misreading. This is the self-grading circularity that the survey of prior LLM verification work identified as its recurring structural weakness, and it is the reason the layer enforces a role separation rather than merely recommending it.

The separation assigns two distinct roles with deliberately opposed objectives. The \emph{elaborator} role is maximally faithful and maximally strict: its objective is to capture the human's intent as precisely as possible, it conducts the Socratic dialogue and owns the pinned oracles, and---critically---it never sees the implementation. The \emph{implementer} role has the narrow objective of satisfying the contract handed to it; it does not negotiate the contract and has no incentive to weaken it, because it never authored it. Because the elaborator is blind to the code, it cannot tailor the oracle to excuse a particular implementation, and because the implementer cannot edit the contract, it cannot relax the oracle to admit a deficient implementation. The verifier, reused unchanged from our prior work as a trusted component, grades the implementer against the elaborator's contract and is authored by neither. The independent authority constraint removes \emph{self-grading} circularity---an agent cannot tailor the oracle to excuse its own code---but it does not by itself decorrelate a shared \emph{misreading of intent} across two roles of the same underlying model, which can encode the same error in both contract and code. That residual is broken not by role separation but by human ratification, whose authority lies outside the model distribution, and is further reduced by drawing the adversarial elaborator from a different model or prompt lineage. The constraint is thus necessary but not sufficient; the provenance tagging introduced next is what makes adherence to it auditable rather than merely asserted.

\subsection{Provenance, Assurance, and Inherited Default Contracts}
\label{subsec:inherited_defaults}

A contract in this pipeline is not a monolithic block of formal text but a \emph{set of obligations}, each carrying two tags. The \emph{provenance} tag records where the obligation came from: \texttt{human-ratified}, \texttt{inherited-default}, \texttt{inferred-unconfirmed}, \texttt{example-pinned}, \texttt{ticket-stated}, \texttt{doc-stated}, or \texttt{code-inferred}. The \emph{assurance} tag records how strongly it has been established against the implementation: \texttt{proved}, \texttt{bounded-checked}, \texttt{tested}, or \texttt{unchecked}. Provenance makes the independent authority constraint auditable, because one can mechanically confirm that every obligation the implementer was graded against carries a provenance other than the implementer's own authorship. Assurance makes the pipeline's ``verified'' claim honest: an obligation that timed out in the solver and fell back to bounded checking is recorded as \texttt{bounded-checked}, never silently promoted to \texttt{proved}. The two tags together turn a contract into a self-describing assurance record rather than an opaque oracle of unknown trustworthiness.

The provenance value \texttt{inherited-default} names the most team- and process-flavored artefact of the layer: \emph{inherited default contracts}. These are organisation- or codebase-level policy obligations that individual elaborations inherit automatically---for example, ``public methods reject \lstinline|null| arguments unless explicitly stated otherwise,'' ``monetary values are never negative,'' or ``returned collections are never \lstinline|null|.'' Such policies encode decisions an organisation has already made once and does not wish to re-adjudicate per method. They sharply reduce the ratification burden, because the elaborator asks the human only about \emph{departures} from policy rather than re-litigating every boundary decision in the catalogue of \cref{subsec:socratic}; a method that conforms to the null-rejection default needs no null question at all. Inherited defaults also make ratification labor scale with the genuinely novel content of each method, which is the practical answer to the objection that example-grounded ratification is too costly to apply across a large codebase. Where a method must depart from a default, the departure is itself surfaced as a behavioural question and, once ratified, retags the affected obligation from \texttt{inherited-default} to \texttt{human-ratified}.

\subsection{The In-Layer Elaboration Loop}
\label{subsec:elaboration_loop}

The mechanisms above combine into a single elaboration loop, given as the numbered procedure below. The loop takes the NL intent and the applicable inherited defaults and terminates with a contract every clause of which is provenance-tagged and either resolved by stated intent, resolved by an inherited default, or ratified through a behavioural example. The adversarial and vacuity checks run before the human is consulted, so the human adjudicates a contract that has already been hardened against the most common weakening failures; the human's questions therefore concern genuine intent decisions rather than drafting defects.

\begin{enumerate}
\item \textbf{Seed.} Obtain a draft contract: for greenfield methods, an LLM few-shot draft from the NL intent; for brownfield methods, the inferrer's recovered draft (\cref{sec:intent_layer} situates this within the broader characterisation phase). Attach inherited default obligations applicable to the method's scope.
\item \textbf{Surface decisions.} Walk the boundary-decision catalogue (\cref{tab:boundary_catalogue}); mark each category as resolved-by-intent, resolved-by-default, or unresolved.
\item \textbf{Adversarial check.} Invoke the independent adversary to seek a behaviour consistent with the NL intent but permitted by the draft contract that a reasonable reading would forbid. If a witness is found, record it as a candidate weakening and add the implied strengthening to the agenda.
\item \textbf{Vacuity check.} Flag any obligation equivalent to \lstinline|ensures true| or guarded by an unsatisfiable precondition; add each flag to the agenda.
\item \textbf{Behavioural ratification.} For every unresolved decision, candidate weakening, and vacuity flag, render the relevant contract behaviour as a concrete input--output example and ask the human in domain terms. Record each ratified answer as a pinned oracle and retag the affected obligation's provenance to \texttt{human-ratified}.
\item \textbf{Reconcile and re-check.} Reformulate the contract to satisfy all pinned oracles and ratified strengthenings. Verify mechanically that every pinned oracle is discharged by the reformulated contract; if any is not, return to step~3 with the violated oracle as a new agenda item.
\item \textbf{Terminate.} When the agenda is empty and all pinned oracles are discharged, emit the provenance-tagged contract for handoff across the pivot to the synthesis and verification layers.
\end{enumerate}

The loop is bounded in practice by the size of the boundary-decision catalogue and the inherited defaults, which together cap the number of questions a single method's elaboration can generate. It is also resumable: because pinned oracles persist independently of the formal text, a later compositional pass that recomputes a caller's obligations re-enters the loop only at step~6, re-checking the reformulated contract against the already-ratified oracles rather than re-interrogating the human.

\subsection{Research Contribution of the Layer}
\label{subsec:intent_contribution}

The named research contribution of the intent elaboration layer is a \emph{ratification protocol that establishes genuine human authority over a formal contract without requiring the human to read formal text}. Its components are individually defensible and jointly novel: behavioural ratification through pinned oracles replaces an unverifiable signature with a re-checkable set of agreed cases; Socratic, catalogue-driven elaboration bounds the human's attention to genuine boundary decisions; adversarial elaboration and vacuity checks, repurposing the inferrer's \lstinline|ensures true| punt~\cite{edmonds_article1} as a weakening alarm, harden the contract against the spec-weakening to which a self-drafting model is prone; the independent authority constraint forbids the implementer from authoring its own oracle; and provenance- and assurance-tagged obligations, refined by inherited organisational defaults, make both the independence of the authority and the honesty of the ``verified'' claim auditable. Together these resolve the ratification gap that prior LLM-based specification work leaves open, in which an inferred or NL-derived contract is treated as ground truth despite never having been genuinely confirmed by the human whose intent it claims to encode. The remaining sections take the contract this layer produces as given and develop the synthesis and verification machinery below the pivot.
